\section{Large-Scale Models and Knowledge Distillation} \label{sec:11-introduction-largescale}
\begin{foldable}
  The history of modern \ac{AI} illustrates the bitter lesson~\cite{sutton2019bitter}: general, compute-driven approaches eventually outperform hand-crafted solutions.
  When AlexNet~\cite{krizhevsky2012imagenet} appeared in 2012, its roughly 60\,M parameters and strong performance on ImageNet signalled that raw model capacity, combined with enough data and compute, could surpass decades of hand-crafted feature engineering.
  An arms race followed that has not slowed since.
  Vision models scaled rapidly: parameter counts of VGGNet~\cite{simonyan2014very} and ResNet~\cite{he2016deep} reached hundreds of millions, while vision transformers soon crossed the 1\,B scale~\cite{dosovitskiy2020image,zhai2022scaling}.
  Language models followed an even steeper trajectory after the Transformer~\cite{vaswani2017attention} architecture arrived, with BERT~\cite{devlin2019bert} featuring 340\,M parameters, followed by GPT-3~\cite{brown2020language} reaching 175\,B parameters the year after.
  As of writing, flagship \ac{AI} models~\cite{anthropic2026assessing,openai2025introducing} are all speculated to have reached trillions of parameters, especially when open-source counterparts too fall in this range~\cite{meta2025llama4,yang2025qwen3}, backed by compute infrastructure demands~\cite{elmeleegy2024demystifying}.
  Their training data have expanded in tandem to tens of trillions of tokens, estimated to have exhausted 80\% of the high-quality English-language surface web~\cite{villalobos2024position}.
  The empirical result is clear: larger models, trained on larger datasets, consistently achieve state-of-the-art performance across vision, language, and multimodal tasks.
\end{foldable}

\begin{foldable}
  Yet scaling success has created a deployment paradox: the models with the strongest benchmark performance are exactly the ones real-world systems struggle most to accommodate.
  Practical deployment is governed not by leaderboard rankings but by hard physical, economic, and legal constraints.
  These constraints operate on two distinct fronts.
\end{foldable}

\begin{foldable}
  The first front concerns the model at inference time.
  Edge devices (smartphones, microcontrollers, or embedded sensors) are subject to fixed limits on silicon area, working memory, power draw, and thermal dissipation.
  A model requiring tens of gigabytes of weights and billions of floating-point operations per forward pass simply cannot execute within these envelopes.
  Safety-critical and real-time applications further compound this with latency requirements.
  For example, an autonomous vehicle processing sensor data cannot tolerate a round-trip to a remote server, a clinical decision-support system must respond within the timescales of patient care, and a fraud detection system must act before a transaction clears.
  Even in cloud settings, where hardware constraints are nominally relaxed, the economics of serving a multi-billion-parameter model at consumer scale are severe.
  Offline deployments such as disaster response, rural healthcare, and maritime or aerospace operations face the same problem in its starkest form: no server is reachable, so the model must run locally or not at all.
  The research community has responded with a family of model compression techniques: pruning, quantization, low-rank adaptation, and \ac{KD}, all aiming to produce smaller, faster models with minimal accuracy loss.
  \ac{KD} in particular transfers knowledge from a teacher model to a student model, requiring no weight access and no structural modification.
\end{foldable}

\begin{foldable}
  The second front concerns the data volume required to build and adapt these models.
  Pre-training on billions of images and documents demands massive infrastructure (storage, pipelines, distributed clusters) available only to a handful of organisations worldwide.
  This incurs expense on training and fine-tuning, often requiring hundreds of hours and terabytes of task-specific data.
  For organisations adapting models frequently, cumulative costs rival the original pre-training expense.
  Energy expenditure compounds this further: a single training run emits carbon dioxide comparable to several passenger vehicles~\cite{strubell2019energy}.
  Beyond these economic pressures, data access is independently constrained by legal and institutional barriers.
  Privacy legislation such as the \ac{GDPR}~\cite{eu2016gdpr} in Europe and the \ac{HIPAA}~\cite{us1996hipaa} in the United States restrict data centralization.
  Datasets are also proprietary, with organisations unwilling to share data that underpins their competitive advantage.
  This creates friction for collaboration and research, since researchers must either forgo real-world data or work through complex legal and ethical requirements to access it.
  Dataset compression techniques, coreset selection, prototype-based methods, and \ac{DD}, address these constraints.
  \ac{DD} in particular optimizes a synthetic set so that models trained on it match full-dataset performance, capturing essential knowledge within a minimal data budget.
\end{foldable}

\begin{foldable}
  Both compression dimensions are independently valuable, but neither is sufficient on its own: a practitioner who compresses only the model still faces the full cost of training data at adaptation time, while one who compresses only the data still deploys a model too large for the target hardware.
  Together, model and dataset distillation\footnote{This thesis uses \textit{knowledge distillation} as an umbrella term for both \textit{model distillation} and \textit{dataset distillation}; the two terms are used explicitly when the contexts are ambiguous.} address the full spectrum of these constraints, in a way that sets them apart from other compression families.
  Both treat compression as a knowledge transfer problem through different mechanisms: model distillation via teacher output, dataset distillation via synthetic data.
  Both are therefore concerned with the fidelity of what is transferred, not just the efficiency of the transfer.
  This concern is what makes the practical gaps non-trivial: when data access is restricted, when the teacher is a black box, or when the target modality is discrete, the central question becomes whether the distilled signal remains faithful to the original.
  The following sections examine where these gaps arise and how the contributions of this thesis address them.
\end{foldable}
